% AI-written, AI cross-reviewed; no human sign-off; not a claim that RH is proved
% RH Notes blueprint source excerpts for Z2.K2
% Verbatim excerpts in the notation of the parent blueprint. The labelled editions and source ranges appear below.
% See the unit YAML for the current proof status. Earlier-round passages retain their original wording.

% Start of atlas edition annotation. The following block is outside the original excerpt.
\paragraph{Status of named inputs in this reading.}
The first entry identifies the local density implication. It assumes the source’s Chebotarev predicate.
\noindent\nolinkurl{GlobalTheory.Published.ChebotarevDegreeOne.dense} --- in-repo constant (\path{Development/AbelianVarietyTheory/PublishedTheorems/GaloisCohomology₂/ChebotarevDegreeOneIffCheb.lean:51--61}); not compared.\par
% End of atlas edition annotation.

\section*{Z2.K2. Cycles on integral matrices and uniform residue evaluation}

% Parent source SHA-256: 1af64922b98ed2310744411b1f18b25cc426c97220646321471836e1a3a409d6
% blueprint/ch_P.tex lines 203 through 308; prop:P-k-few
\begin{proposition}[Separation by evaluations]\label{prop:P-k-few}
\lean{LateralTheory.Q.SFewOfCut.sFewOf}
\uses{prop:P-k-layer}
Set
\[
 b=2E+(2E+1)(c+2E+1),\quad N_0=4E+6,\qquad
 T_i=\left|H_i(\mathrm{GL}_{N_0}(\mathcal O_M),\mathbf Z)_{\rm tors}\right|.
\]
If a finite set $S$ of eligible $\alpha$ has pairwise distinct
$N_a(\alpha)^{\ell^b}$ in $F^\times/(F^\times)^q$, then
$|S|\le q^r T_{2E+1}T_{2E}$. The torsion here is the entire integral
torsion group (\nolinkurl{torsHGL}, \path{GeneratorCountTuples/TorsBoundPow.lean:13}).
\end{proposition}
\begin{proof}
We specify the chain and detection inputs at the point they enter.

\emph{K2, cycle construction and evaluation.} Section~\ref{sec:P-r3-K247} constructs the relation and interpolating multiplicities. The scalar evaluator, integral transport and uniform reduction are proved in the K2 argument at \ref{sec:P-r4-K2}.
Put $m=2E+1$. For each eligible $\alpha$ with divisor $u$ there
exist a degree-$m$ cycle $y_\alpha$ modulo $q$ in
$\mathrm{GL}_{N_\alpha}(\mathcal O_M)$ and a nonzero
$b_\alpha\in\mathcal O_F$. For every reduction $\mathcal O_F\to k$ to a
finite field that is surjective on $\mathcal O_M$ and preserves
$b_\alpha\ne0$, its evaluation by any admissible additive functional
with index character $v:k^\times\to\mathbf Z/q$ is
\[
  P_E\,v\big(\overline{N_a(\alpha)}\big),\qquad
  P_E=\prod_{\substack{0\le i<2E+1\\ i\ne E}}
                   (\kappa^E-\kappa^i).
\]
Here the bar means a residue represented by an integral numerator and
denominator with nonzero denominator. At a prime away from $\ell$ the
root $\zeta$ stays primitive. For a chosen $v$ one can choose such a
functional common to all the $y_\alpha$; its pullback is a family
vanishing on boundaries, stable under enlargement of the matrix size and
additive under block repetition on cycles. These are exactly the uses of
\nolinkurl{SClass}, \nolinkurl{STheta} and \nolinkurl{SPhi}
(\path{LemmaKOfSteps.lean:54--69,84--96}), with their proofs instantiated in
\nolinkurl{lemmaK} (\path{ClassCounting/LemmaKOfLitN.lean:28--31}).
The degree is $2E+1$ because the input there is indexed by $e=E-1$.
In particular it includes $E=1$.

Exclude primes dividing $\ell\prod_{\alpha\in S}b_\alpha$; the product
is nonzero even for the empty set. A primitive $\ell^\nu$-th root remains
primitive at each remaining prime: if its $\ell^{\nu-1}$-st power
reduced to $1$, the geometric sum for a primitive $\ell$-th root would
give $\ell=0$ in the residue field. Its multiplicative group is cyclic of
order divisible by $q$. Taking an exponent modulo $q$ in a cyclic
generator gives an index character $v$ with kernel the $q$-th powers
(\nolinkurl{prim_map}, \nolinkurl{exists_index},
\path{SFewOfCut.lean:85--117}).

Use all the allowed primes $P$ at which $\mathcal O_M\to k(P)$ is
surjective. Assign to $\alpha$ the vector of evaluations of
$\ell^{2E}y_\alpha$ in $\prod_P\mathbf Z/q$. If two vectors agree, then
$\ell^{2E}P_E$ kills the difference of the index characters of the two
twisted norms at every such $P$. In $\mathbf Z/q$ the multiple
$\ell^{(2E+1)(c+2E+1)}$ belongs to the ideal generated by $P_E$.
For this estimate each factor is a unit times
$\kappa^{|E-i|}-1$, with $1\le |E-i|<2E+1<\ell^c$;
the $\ell$-adic valuation is $c+v_\ell(|E-i|)$ (at $2$, $c\ge2$).
There are at most $2E+1$ factors. Removing the prime-to-$\ell$ part
modulo $q$ gives the asserted ideal containment even when the bound
already exceeds $\nu$ (\nolinkurl{pe_dvd}, \path{SFewOfCut.lean:67--83}).
The ratio of the two $\ell^b$-powered norms consequently reduces to a
$q$-th power at every allowed prime. To pass from residues to $F$, choose a $q$-th root $x$ of this ratio
$\beta$ in an algebraic closure. For $g\in G_F$ the set of $\sigma$ with
$\sigma x=gx$ is a nonempty open coset. The density input \emph{K3, proved in Proposition~\ref{prop:P-k-cheb}} says that the Frobenius elements with degree-one prime in $F$
over $\mathbf Q$, allowing the choice of a prolongation, meet every such
open subset even after excluding finitely many primes
(\nolinkurl{ChebotarevDegreeOne.dense}, \path{SDetectCut.lean:111--126}).
The residue field at the chosen prime has prime order. Its prime subfield
already fills it, so $\mathcal O_M$ surjects onto it. By the test, $\beta$
reduces to some unit $\bar t^q$. Write $t$ for an integral representative.
The factorization of $x^q-t^q$ over $\mu_q$ gives
$x\equiv\zeta^i t$ at the chosen prolongation for some $i$. Indeed,
the product of the linear factors has positive additive valuation, so
at least one factor does. Thus $x$ is a unit with residue in the base
residue field. Its Frobenius conjugate $\sigma x$ has the same residue,
and $\sigma x/x$ is a $q$-th root of unity reducing to $1$. Away from
$q$ that root must equal $1$: for a nontrivial root the geometric sum
of length $q$ is zero and its reduction would give $q=0$.
Consequently $gx=\sigma x=x$ for every $g$. The fixed-field theorem gives
$x\in F^\times$ and $\beta=x^q$. This proves the implication
\nolinkurl{SDetect} (\path{LemmaKOfSteps.lean:109--115}) from the
degree-one density assertion (\nolinkurl{frob_fix}, \nolinkurl{goal},
\path{SDetectCut.lean:43--99,101--142}). Distinctness now forces the two
members of $S$ to coincide.

Finally the vectors form evaluations of a \emph{single} admissible
family with values in the product module. Its operators, and the
identities for boundaries, stabilization and repeated blocks, are
defined coordinate by coordinate. The required bound on these images, K4, is as follows: for even $k$
with $k+1<r$, any such family with values in a $\mathbf Z/q$-module and
any finite set of values $\ell^k\Phi(y)$ on degree-$(k+1)$ cycles in
arbitrary matrix sizes, its cardinality is at most
$q^r$ times the torsion orders in degrees $k+1,k$ in rank
$2(k+1)+4$. See \nolinkurl{SUpper}
(\path{LemmaKOfSteps.lean:100--107}) and \nolinkurl{sUpperOf}
(\path{SUpperOf.lean:21--66}), whose reduction uses the Bockstein and
the rank bound \nolinkurl{BOdd} (\path{LemmaKOfSteps.lean:11--15}).
All matrix sizes occurring in a finite set can first be stabilized into
one common rank. With $k=2E$ the stable rank is $N_0$. Injectivity just
proved gives the stated inequality (\nolinkurl{sFewOf},
\path{SFewOfCut.lean:153--211}).
Proposition~\ref{prop:P-k-bock} derives K4 from a stable torsion comparison and a rank bound for the indicated classes. \emph{Status: complete modulo K4.}
\end{proof}

% blueprint/ch_P.tex lines 1218 through 1419; sec:P-r4-K2
\paragraph{K2. The scalar cochain and the cycle on the integral matrix group.}\label{sec:P-r4-K2}
Retain $m=2E+1$, $E=e+1$, $R=\mathbf Z/q$, $a=\kappa^E$ and the
notations of Proposition~\ref{prop:P-k-few}. First write
$z=\sum_{j_1,\ldots,j_E\in C}[j_1|1|\cdots|j_E|1]$.
In its bar boundary the first face sums to zero because $q=0$ in $R$.
Successive faces numbered $1,2$ through $2E-1,2E$ have equal sums:
merging a $1$ with its neighbor translates the corresponding freely
summed $j_s$ by $1$, and the last face deletes the last $1$. Opposite
signs cancel. For a group element $g$, write $P_g$ for the signed sum
of insertions of $g$ in a word, conjugating each subsequent letter by
$g^{-1}$. Its boundary plus $P_g\partial$ is conjugation by $g^{-1}$
minus the identity: the two faces meeting an inserted letter give the
same word in adjacent insertions, so only the two end faces survive.
With the boundary convention used here, the double prism has boundary $P_{gh}-P_g-P_h(\operatorname{conj}_{g^{-1}})_*$ on cycles, after choosing the negative of the projected double insertion as primitive. The signed calculation follows below.
Homomorphisms commute with each face and insertion. For the trivial
group the chain group has one basis vector in every degree and the
positive-degree differentials alternate between $0$ and $1$.
See \nolinkurl{zc}, \nolinkurl{sZc}, \nolinkurl{sChain},
\nolinkurl{sPrism}, \nolinkurl{sPrismMul}
(\path{ClassCounting/HasVal.lean:17--42},
\path{ClassCounting/FirstFace.lean:66}, \path{ClassCounting/SChain.lean:110},
\path{ClassCounting/SPrism.lean:6}, \path{ClassCounting/SPrismMul.lean:134--153,194}).

\emph{Scalar evaluation.} Fix $k$, a primitive $q$th root $\xi$ in
$k^\times$ and a homomorphism $v:k^\times\to R$
(written additively on values). Choose a generator $d$ of $k^\times$,
write $|k^\times|=qr$ and $\xi=d^{ru}$ with $(u,q)=1$.
Choose $s$ with $su\equiv1\pmod q$ and integral lifts $L(g)$ of
$s\log_d g\bmod qr$. Then $L(\xi^x)\equiv rx\pmod{qr}$.
The integer $(L(g)+L(h)-L(gh))/(qr)$ defines $U(g,h)\in R$.
It is symmetric and satisfies
$U(h,j)-U(gh,j)+U(g,hj)-U(g,h)=0$, as expansion of the four
numerators shows. On a word of length $m$, put
\[
 t([g_1|\cdots|g_m])=
 \left(\prod_{i=1}^E U(g_{2i-1},g_{2i})\right)v(g_m).
\]
The bar coboundary is zero by the cup-product rule applied to $U$
and $v$ (equivalently, group the first three faces and use the displayed
identity successively). In $t(P_g w)$, the two insertions before and
after the first letter cancel by symmetry of $U$. Factor out the first
pair and repeat. Only the last insertion contributes, with sign $+1$.
For $\chi_x(j)=\xi^{xj}$, summing the integral numerators of
$U(\chi_x(j),\xi^x)$ over $j\in C$ telescopes to $qL(\xi^x)$,
since $j\mapsto j+1$ permutes $C$. Their quotients sum to
$L(\xi^x)/r\equiv x\pmod q$. The $E$ sums in $z$ are independent,
so $t(P_g(\chi_x)_*z)=x^E v(g)$.
This gives \nolinkurl{sCoc} (\path{ClassCounting/FProd.lean:11--96,101--150}).

\emph{Transfer to finite-dimensional spaces.} For $V$ over $k$ choose
$\lambda(v)\in k^\times$ on nonzero vectors with
$\lambda(av)=a\lambda(v)$, by choosing one vector on every line.
For a representation $\rho$ put $xg=\rho(g)^{-1}x$ for projective
lines $x=[w]$ and
$c(g,x)=\lambda(w)/\lambda(\rho(g)^{-1}w)$.
This is independent of $w$ and
$c(gh,x)=c(g,x)c(h,xg)$. Transfer a bar word by summing over all
starting lines $x_0$ the words
$[c(g_1,x_0)|\cdots|c(g_m,x_{m-1})]$, with $x_i=x_{i-1}g_i$.
The middle faces commute with transfer by the product identity, and
the first face commutes because $x_0\mapsto x_1$ is a permutation.
The last face simply drops the last label. Thus composition with $t$
kills boundaries. For a second choice $\lambda^{\prime}$, put $b([w])=\lambda(w)/\lambda^{\prime}(w)$. Then $c(g,x)b(xg)=b(x)c^{\prime}(g,x)$. The alternating insertions of the vertical edges give $dQ+Qd=T_{c^{\prime}}-T_c$, as proved below.
The value on cycles therefore survives a change of choices and an
isomorphism of representations. In dimension one there is a single
line and $c(g,x)$ is the scalar itself
(\nolinkurl{tr_cob}, \nolinkurl{theta_iso}, \nolinkurl{theta_one},
\path{ClassCounting/STrans.lean:16--25,82--95,130--165,169--249,274--294}).

For $V\oplus W$ use $\lambda_V(v)$ when $v\ne0$ and $\lambda_W(w)$
otherwise. A line of the first kind projects to a line of $V$; after
fixing its $V$-representative its lifts are parametrized by $W$.
Lines of the second kind are exactly the lines of $W$. For a block
diagonal action every label is the label of the projected path.
The multiplicities are $|W|=|k|^{\dim W}=1$ in $R$ and $1$,
respectively. This gives direct-sum additivity, including zero
dimensions. The construction of the family $\theta$ uses only
$k,\xi,v$. Thus one family evaluates all the classes on integral matrix groups, with coefficients in $R$, at that reduction. Together with the determinant and repetition
arguments above this gives the required \nolinkurl{sTrans},
\nolinkurl{sThetaOf} and \nolinkurl{sPhi}
(\path{ClassCounting/STrans.lean:303--404,411--413},
\path{ClassCounting/SThetaOf.lean:192--273},
\path{ClassCounting/PiRepApply.lean:82--154}).

\emph{Transport with prescribed determinants.} Let $D$ be the
eligibility divisor of $\alpha$, $J=C^2$, $\Lambda_0=\mathcal O_F^J$,
and $\Lambda_D=\bigoplus_{j\in J}I_j$ where $I_j$ has orders
$\eta(j)D$. Give the $j$th coordinate of $F^J$ the character
$(s,t)\mapsto\zeta^{j_1s+j_2t}$, defining $\rho$.
Along $d_t$ its weight is $p_t(j)$. For $t\ne t_0$, at every
weight the orders of the product of the $I_j$ in that weight sum to
zero by the interpolation relation. The Steinitz argument above gives
a transport from $\Lambda_0$ to $\Lambda_D$ with determinant $1$
in each weight. For $t=t_0$ set
\[
 p(Y)=\prod_{0\le i<m,\ i\ne E}(Y-\kappa^i)=\sum B_iY^i,
 \qquad \beta=\prod_i(\gamma^{-i}\alpha)^{B_i}.
\]
Writing $(Y-a)p(Y)=\sum A_iY^i$ gives
$\operatorname{ord}_P\beta=\sum A_iD(\gamma^iP)$, by substituting
$\operatorname{ord}_{\gamma^iP}\alpha=
D(\gamma^{i+1}P)-aD(\gamma^iP)$.
See \nolinkurl{cGtNe}, \nolinkurl{cBetaDiv}, \nolinkurl{polAct}
(\path{LatticeBlockClasses/CGtNe.lean:7},
\path{LatticeBlockClasses/CBetaDiv.lean:33},
\path{LatticeBlockClasses/DiagApply.lean:42--52,74--86}).

On a coordinate of weight $p_{t_0}(j)=\kappa^i$, change coordinates
by $s_j=\gamma^{-i}$, and use $s_j=1$ off this orbit. This choice is
well defined: $v_\ell((1+\ell^c)^n-1)=c+v_\ell(n)$ for $n>0$,
by binomial expansion first for $n$ prime to $\ell$ and then for each
successive $\ell$th power ($c\ge2$ when $\ell=2$). Hence the order
of $\kappa\bmod q$ equals the specified $|\gamma|=\ell^{\nu-c}$.
All orbit weights now have exponent $1$
(the character $j'\mapsto\zeta^{j'}$). Other weights stay unchanged.
Their total ideal divisor is zero. At the new weight $1$ the divisor is
\[
 \sum_{j:\ p_{t_0}(j)\in\langle\kappa\rangle}
 \eta(j)D(s_j^{-1}P)=\sum_i A_iD(\gamma^iP)
 =\operatorname{ord}_P\beta,
\]
by the $p_{t_0}$-image of $\eta$; coincident powers contribute the
sum of their coefficients. The weight-$1$ block contains $(1,0)$.
Steinitz on the twisted blocks therefore gives matrices $A_x$ with
determinants $\delta_x=\beta$ for $x=1$ and $1$ elsewhere (empty
blocks have determinant $1$). Conjugating back by the coordinate
change produces an $M$-linear transport $g_{t_0}$ commuting with
$\rho d_{t_0}$, since each twisted character is scalar on its block.
Off $t_0$ use the same notation with $s_j=1$ and all $\delta_x=1$.
This proves the precise \nolinkurl{Carries} assertion
(\nolinkurl{cGtEq}, \nolinkurl{cBlock},
\path{LatticeBlockClasses/GtEqZeta.lean:11--52,54--106},
\path{LatticeBlockClasses/CBlock.lean:83},
\path{LatticeBlockClasses/CJoin.lean:12--37}).

\emph{One localization for every reduction.} For
$\Gamma=\operatorname{Gal}(F/M)$ the map
$F\otimes_M F\to\prod_{\tau\in\Gamma}F$ taking
$u\otimes v$ to $(u\tau(v))_\tau$ is an isomorphism: a primitive
element has distinct conjugates and the Chinese remainder theorem
splits its minimal polynomial. The same map on integral lattices is
an isomorphism after inverting a single $b_0\in\mathcal O_M\setminus\{0\}$
since its kernel and cokernel are finite torsion $\mathcal O_M$-modules.
Choose $b\in\mathcal O_M\setminus\{0\}$ divisible by $b_0$ and by
an element below every $P$ in the support of $D$ (one may multiply
by their underlying rational primes). After inverting $b$ the lattices
$\Lambda_0$ and $\Lambda_D$ agree, so all the $g_t$ and their inverses
act on the same localized lattice. Since $b$ is fixed by $\Gamma$,
all conjugates of $\alpha$ and $\beta$ are units there by the divisor
identities above. For any allowed $r:\mathcal O_F\to k$, tensor the
integral splitting on its first factor by $r$. This identifies
$k\otimes_{\mathcal O_M}\mathcal O_F^J$ with $k^{\Gamma\times J}$
by coordinates $r\tau$. Surjectivity on $\mathcal O_M$ makes this a
quotient of $\Lambda_0$ itself. The action on the localization
therefore reduces functorially at this $r$; its restriction to the
integral stabilizer is the reduction of the original lattice action.
A finite torsion-free module over the Dedekind ring $\mathcal O_M$ is
projective, so choose a complement to $\Lambda_0$ in a finite free
module and give it the trivial action. This choice precedes $r$. Its
positive-degree homology contribution is zero by the bar calculation
above. These are the splitting and freeness assertions of
\nolinkurl{cSplit}, \nolinkurl{cLatLoc}, \nolinkurl{cFree} and
\nolinkurl{cTriv} (\path{LatticeBlockClasses/CSplit.lean:5--44,47--92},
\path{LatticeBlockClasses/CLatLoc.lean:87},
\path{LatticeBlockClasses/RedApplyMap.lean:65},
\path{LatticeBlockClasses/CTriv.lean:6},
\path{LatticeBlockClasses/CJoin.lean:87--136}).

Write $\omega(\tau)$ for the exponent of $\tau$ on $\zeta$ modulo
$q$. In the split coordinates the $x$-block at $\tau$ uses the entries
$(\tau s_j,j)$ for all $j$ of twisted weight $x$. This is a
permutation of the coordinates. The character there has exponent
$\omega(\tau)x$ and the transport matrix is $r(\tau A_x)$, acting
entrywise, with determinant $r(\tau\delta_x)$.
Indeed $s_j((g_tw)_j)=(A_x(s_iw_i)_i)_j$ and applying $r\tau$ gives
exactly these entries. Each $A_x$ and its inverse preserve the standard
localized block since its source and target agree there.
The determinant evaluation on a scalar block and the boundary
calculation for the cycle with coefficients in $R$ now give, after scaling by
$n_{t_0}^{-1}$, exactly $\sum_\tau\omega(\tau)^E v(r(\tau\beta))$.
To compute it put $S(\alpha)=\sum_\tau\omega(\tau)^Ev(r(\tau\alpha))$.
Reindexing $\tau$ by $\tau\gamma^i$ gives
$S(\gamma^{-i}\alpha)=a^i S(\alpha)$ in $R$.
Additivity of $v$ on units thus gives
\[
 S(\beta)=p(a)S(\alpha)
 =P_E\,v\!\left(r\!\left(\prod_{j=0}^{|\Gamma|-1}
                       (\gamma^j\alpha)^{a^j}\right)\right).
\]
All residues here use the one localization; clearing a power of $b$
provides the numerator and denominator in \nolinkurl{Res}. The
calculation for an already admissible family uses only its scalar
block law, so it also applies whenever the reduced root has $q$th
power $1$ and such a family is given. Primitivity enters the separate
existence construction of that family.
This finishes K2 for $E\ge1$, $2E+1<\ell^c$ and the stated layer
and eligibility hypotheses
(\nolinkurl{cVal}, \nolinkurl{cBetaVal}, \nolinkurl{cJoin},
\nolinkurl{sClassOf}, \path{LatticeBlockClasses/BlkEquiv.lean:9--75},
\path{LatticeBlockClasses/CBetaVal.lean:31--119},
\path{LatticeBlockClasses/CJoin.lean:131--187},
\path{ClassCounting/SClassOf.lean:21--23}).

